%!TeX root=./tese.tex
% Baselines chapter. Every number is ledgered in progress/results-ledger.md:
% "Static-optimal grid, final", "3.1 close-out checks", "3.2 reactive LLM
% baseline", "4.4 rider (c) reactive target ladder", "4.4 fleet matrix"
% (reactive arm). Artifacts: simulation/results/grid/, reactive32/,
% riders44/ladder_*, results_44.jsonl.
\chapter{Provisioning Baselines: the Static Optimum and Reactive Control}
\label{cap:baselines}
\chaptermark{Provisioning Baselines}

Any claim that predictive provisioning pays must be made against the
two policies a practitioner would reach for: the best \emph{static}
fleet and a \emph{reactive} autoscaler. Both are established here on the
platform of Chapter~\ref{cap:twopool}, under the evaluation SLA
defined there (per window: language TTFT p99 $<$ \SI{1}{\second} and ITL
p99 $<$ \SI{100}{\milli\second}; encoder delay p99 $<$
\SI{1}{\second}), and scored on provisioned GPU-time \emph{per hardware
class} (RTX~5090-class worker-seconds for the encoder pool, H100
GPU-hours for the language pool; never summed across classes).

\paragraph{From order or from readiness: the pools bill differently.}
The two classes bill from different instants. Language-pool GPU-hours
bill from instance \emph{order}, not readiness: the discrete-event
mocker integrates the \emph{provisioned} worker count --- active,
starting-up and draining alike --- so the startup ramp and scale-down
drain tail are already inside every reported language GPU-hour, and
\emph{no cold-start surcharge may be added to it}. The encoder pool bills
from \emph{readiness}, its cold-start reported separately as a warm-up
column (Section~\ref{sec:cl-ledger-enc}). Mixing the two conventions is a
second reason, beyond hardware class, the classes are never summed; the
order-time convention is read off the mocker's collector and corroborated
over the language runs (Appendix~\ref{ap:artifacts}).

Both baselines run on the same windows and actuation mechanics the
predictive policy of Chapter~\ref{cap:closedloop} will have. Cold-start
leads are matched on the encoder side; the fourteen headline language
reactive configurations, however, ran with \emph{instant} workers (no
startup lead) --- an optimistic zero-startup lower bound for reactive
control, so their tail-SLA failures below are, if anything, conservative
(Section~\ref{sec:base-reactive-llm}).

\section{Protocol}
\label{sec:base-protocol}

The workload windows are the six largest burst episodes of
Chapter~\ref{cap:trace-eda} (onsets 1--404 minutes) plus the calmest full
day of the trace (2024-10-19; peak \SI{57}{img/s}) as the steady
reference --- all shaded on the week in Figure~\ref{fig:eda-windows}. Each burst window
brackets its episode with a $\pm$10-minute calm margin (the median-crossing
envelope of Section~\ref{sec:eda-vision}), so every replay begins and ends
with a drained encoder queue --- arrivals may end with the queue still
nonempty, but the replay runs on until it drains and that tail is inside
the replay span --- and this padded span sets the
window's GPU-time denominator. Two episodes
(burst2/6) fall in the trace's first 1.4 days, reserved as forecast
context (temporal split, Chapter~\ref{cap:closedloop}); they are graded
here only by the forecast-free policies --- the static grid and both
reactive baselines (ARIMA loop internal, not the archived forecast these
windows lack). Each policy produces, per window, a worker schedule for
the encoder pool ($C(t)$ workers) and a configuration or replica schedule
for the language pool ($\mathrm{TP}\times N$ H100 workers), graded on the
two-pool replay. Cold-start lead is modeled explicitly: a scale-up
decision issued at $t$ delivers usable capacity at $t+L$. The encoder
pool is swept over $L \in \{15,30,45,300\}$\,\si{\second}, above the
measured 3--5\,\si{\second} model-side floor. On the language pool,
leads $L \in \{60,\ldots,300\}$\,\si{\second} (the weight-loading
regime) apply to the predictive runs and to the burst1 reactive
sensitivity below; the fourteen headline language reactive runs model
no startup lead (Section~\ref{sec:base-reactive-llm}).
Scale-downs take effect immediately, guarded by a 3-minute hysteresis.

\section{The static optimum: burst-shaped, calm-day expensive}
\sectionmark{The static optimum}
\label{sec:base-static}

A grid search (encoder workers $C$ crossed with language shapes
$\mathrm{TP}\times N$, refined around the feasibility frontier; 101 points
archived, Appendix~\ref{ap:artifacts}), plus a superseding burst5
stability test that corrected its raw grid endpoint ($C^{*}$ 16 $\to$ 8;
the original divergence flag was an endpoint artifact), yields the
effective optima per window (Table~\ref{tab:base-grid}).

\begin{table}[htbp]
\centering
\caption{Static-optimal fleet per window under the evaluation SLA.
$C^{*}$: minimum feasible encoder-worker count; LLM shape: cheapest
feasible $\mathrm{TP}\times N$. GPU-time per hardware class over the
window span.}
\label{tab:base-grid}
\small
\begin{tabular}{@{}lrrlrr@{}}
\toprule
Window & Span (h) & $C^{*}$ & LLM shape & 5090 GPU-h & H100 GPU-h \\
\midrule
burst1 & 8.32 & 6 & TP2$\times$2 & 49.9 & 33.3 \\
burst2 & 1.90 & 8 & TP2$\times$2 & 15.2 & 7.6 \\
burst3 & 0.45 & 7 & TP2$\times$2 & 3.2 & 1.8 \\
burst4 & 3.15 & 7 & TP2$\times$2 & 22.1 & 12.6 \\
burst5 & 1.35 & 8 & TP2$\times$2 & 10.8 & 5.4 \\
burst6 & 1.37 & 6 & TP2$\times$2 & 8.2 & 5.5 \\
steady (24 h) & 24.0 & 2 & TP1$\times$2 & 48.0 & 48.0 \\
\bottomrule
\end{tabular}
\end{table}

The encoder $C^{*}$ minimizes the all-request encoder-delay p99
(Chapter~\ref{cap:twopool}); under the stage-conditional metric the
burst optima rise one to two workers but the deployable static $C{=}8$
is unchanged, so the metric reshapes only the per-window table, not the
fleet behind Chapter~\ref{cap:closedloop}'s cost headline.

Three structural facts follow. First, \textbf{every burst window forces
the language pool to TP2, and the gain is the parallelism layout, not
the GPU budget}. At TP1 no replica count is feasible: the burst's
${\approx}36$k-token frames-class prefills pin TTFT p99 at
${\approx}$\SI{1285}{\milli\second} even at $N{=}16$, a \emph{prefill
floor} replica scaling cannot buy down. Four TP1 replicas and one
TP2$\times$2 group are both four H100s, yet TP1$\times$4 fails all five
tabulated burst windows while TP2$\times$2 passes each --- a
1.55--1.74$\times$
tail cut from layout alone (Table~\ref{tab:base-tp}): intra-request
tensor parallelism shortens the single long prefill that independent
replicas can only queue. Second, the encoder requirement
swings from 2 workers on the calm day to 8 on burst5: a fourfold
fleet-sizing swing, the operational consequence of the
order-of-magnitude workload-rate mismatch of Chapter~\ref{cap:trace-eda}. Third, the single static fleet that survives
\emph{all} windows, $C{=}8$ encoders plus TP2$\times$2 (4 H100s),
hereafter the \emph{deployable static}, overprovisions the calm
day by $4\times$ on the vision side and $2\times$ on the language side.
It also carries a fleet-\emph{shape} difference (TP2 versus TP1) that no
replica autoscaler can morph at runtime. Close-out checks confirm the
optima are robust to the router's tuning knobs (35 corner re-runs;
no frontier movement).

\begin{table}[htbp]
\centering
\caption{Iso-cost tensor parallelism on the language pool: TP1$\times$4
and TP2$\times$2 both cost four H100s. TTFT p99 per burst window under the
\SI{1000}{\milli\second} SLO (ITL p99 within bound in both columns); AIC
timing on each window's composed trace at its archived grid encoder
count --- burst1--4 at the window optimum, burst5 archived at $C{=}16$
(predating its corrected $C^{*}{=}8$).}
\label{tab:base-tp}
\small
\begin{tabular}{@{}lrrr@{}}
\toprule
Window & TP1$\times$4 p99 (ms) & TP2$\times$2 p99 (ms) & Tail cut \\
\midrule
burst1 & 1{,}318 \emph{(fail)} & 763 \emph{(pass)} & 1.73$\times$ \\
burst2 & 1{,}331 \emph{(fail)} & 766 \emph{(pass)} & 1.74$\times$ \\
burst3 & 1{,}315 \emph{(fail)} & 846 \emph{(pass)} & 1.55$\times$ \\
burst4 & 1{,}317 \emph{(fail)} & 792 \emph{(pass)} & 1.66$\times$ \\
burst5 & 1{,}316 \emph{(fail)} & 781 \emph{(pass)} & 1.68$\times$ \\
\bottomrule
\end{tabular}
\end{table}

\section{Reactive control on the language pool}
\label{sec:base-reactive-llm}

The reactive baseline is Dynamo's SLA planner, the production controller
of the platform this thesis targets, run through the same replay path in
both its shipped load-based variant (A) and its SLA-target variant (B,
ARIMA-driven throughput loop plus load loop). Across the seven windows
this yields 14 configurations, each at its window's optimal TP shape.
These headline runs model no startup lead --- scale-ups deliver workers
instantly, an optimistic zero-startup lower bound for reactive control.
\textbf{All fourteen configurations satisfy the planner's own control
law (window-mean TTFT/ITL under its 500/50\,ms targets) while violating
the tail SLA} (TTFT p99 up to 3{,}495\,\si{\milli\second}; burst windows
additionally fail ITL p99). Instant scale-up favors the reactive
baseline, so the 14/14 tail failure is conservative. Variants A and B
fail alike, so loop tuning is not the issue. The defect is the \emph{mean-based capacity
conversion}: sizing replicas so mean throughput meets mean demand leaves
no room for the tail the SLA prices.

Nor can spend buy the tail back under realistic lead. Holding a 60-second
replica lead (this sensitivity, unlike the headline runs, does model the
startup lead) and tightening the planner's window-mean targets from
400/40 to 200/20\,\si{\milli\second} on burst1 raises GPU-hours
monotonically (18.6\,$\to$\,20.4) while TTFT p99 stays pinned near
1{,}616\,\si{\milli\second} --- a 10\% spend increase buying nothing.
Burst1's 404-minute onset gives a trailing window ample time, yet
tightening the mean targets still cannot buy the tail: the failure is the
conversion, not thrift or timing (the complement of the fast-onset
encoder case below).

\section{Reactive control on the encoder pool}
\label{sec:base-reactive-enc}

The encoder pool has no native reactive controller, so the fair strawman
is \emph{identically informed} with the next chapter's predictive policy:
the same worker-mapping law (workers $= \lceil \text{demand} / \rho
\rceil$ at per-worker capacity rating $\rho$ \si{img/s}, hysteresis 3
minutes, lead $L$), fed by the trailing window of the \emph{observed
image rate} --- the correct per-pool signal, observed not forecast.
Table~\ref{tab:base-react-enc} grades it at the operating rating
($\rho{=}15$) across cold-start leads.

\begin{table}[htbp]
\centering
\caption{Reactive encoder scaling on the observed image signal
($\rho{=}15$): encoder-delay p99 in ms by cold-start lead $L$; bold =
SLA violation. $L^{*}$ = largest \emph{tested} lead at which the SLA
holds.}
\label{tab:base-react-enc}
\small
\begin{tabular}{@{}lrrl@{}}
\toprule
Window & $L{=}15$\,s & 300\,s & $L^{*}$ \\
\midrule
burst1 (onset 404\,min) & 588 & 638 & 300\,s \\
burst2 (onset 48\,min) & 588 & \textbf{4{,}282} & 45\,s \\
burst3 (onset 1\,min) & \textbf{9{,}481} & \textbf{54{,}170} & none \\
burst4 (onset 60\,min) & 588 & \textbf{3{,}063} & 165\,s$^{\dagger}$ \\
burst5 (onset 35\,min) & 588 & \textbf{6{,}729} & 45\,s \\
burst6 (onset 53\,min) & 588 & \textbf{4{,}606} & 45\,s \\
steady (24\,h) & 548 & 884 & 300\,s \\
\bottomrule
\end{tabular}

\vspace{2pt}
{\footnotesize $^{\dagger}$ A supplementary sub-grid cell outside the
declared $\{15,30,45,300\}$\,s factorial: burst4 passes at
\SI{165}{\second} (\SI{904.7}{\milli\second}) and fails at the next
tested lead, \SI{300}{\second}; the 165--300\,s interval is unswept, so
\SI{165}{\second} is the largest tested passing lead, not an identified
feasibility threshold. The other medium onsets (burst2, burst5, burst6)
are limited to \SI{45}{\second}.}
\end{table}

The pattern tracks onset speed. Slow ramps
(burst1, the calm day) are reactive-scalable at any tested lead; medium
onsets (35--60 minutes) survive only short leads (largest tested passing
leads 45--\SI{165}{\second}; three of the four windows are limited to
\SI{45}{\second}, below what pod-scheduled GPU workers realistically
offer, though burst4 tolerates a 2.75-minute lead). The one fast onset,
burst3, ramps baseline-to-peak in a minute and defeats the controller at
\emph{every} lead, including \SI{15}{\second} (p99 encoder delay
9.5--54\,\si{\second}) --- while burst6, at a comparable peak
(\SI{337}{} vs burst3's \SI{355}{img/s}) and essentially the same mean
image rate, request rate, and per-request image density, is caught at
short leads ($L^{*}{=}$\SI{45}{\second}), an outcome consistent with its
much slower merged onset (53 minutes rather than one). So
watching the right signal is necessary but not sufficient: against a fast
onset no affordable lead delivers workers before the queue does the damage.

\section{Summary}
\label{sec:base-summary}

The static optimum is burst-shaped: deployable only as 8 encoders +
TP2$\times$2, i.e.\ $4\times$/$2\times$ calm-day overprovisioning plus a
TP-shape change no replica scaler performs (RQ1's missizing, priced).
The reactive alternative fails on both pools for one structural reason
--- information arriving too late and converted through means: the
fourteen language-pool configurations meet their own window-mean law
while violating the tail SLA even with instant workers,
spend cannot buy the tail back under lead,
and even the correct encoder signal read reactively loses every window
whose onset is faster than the lead. At full-trace scale
(\num{1000000} requests), the deployable $C{=}8$ fleet passes, the
calm-day $C{=}2$ fleet fails, and reactive encoder control holds the
bound for 0.687$\times$ the deployable fleet's worker-time --- a real
aggregate-scale saving, but graded in aggregate only (the masking of
Section~\ref{sec:cl-ledger}), so it does not overturn the per-window
verdicts: reactive control still loses burst3 graded on its own. Whether
\emph{forecasting} those same per-pool signals closes
the gap is the subject of the next chapter.
